% Generated draft: requires mathematical and source review.
\documentclass[11pt]{article}
\usepackage[margin=25mm]{geometry}
\usepackage{fontspec}
\setmainfont{Latin Modern Roman}
\usepackage{amsmath,amssymb,mathtools,unicode-math}
\setmathfont{Latin Modern Math}
\usepackage{enumitem,needspace,xurl,hyperref}
\setlength{\parindent}{0pt}
\setlength{\parskip}{5pt}
\setlength{\emergencystretch}{3em}
\newcommand{\N}{\mathbb{N}}
\newcommand{\Z}{\mathbb{Z}}
\newcommand{\Q}{\mathbb{Q}}
\newcommand{\R}{\mathbb{R}}
\newcommand{\C}{\mathbb{C}}
\newcommand{\F}{\mathbb{F}}
\newcommand{\A}{\mathbb{A}}
\newcommand{\PP}{\mathbb{P}}
\newcommand{\E}{\mathbb{E}}
\newcommand{\eps}{\varepsilon}
\newcommand{\dd}{\,\mathrm{d}}
\begin{document}
\section*{UnsolvedMath ID: 2303020}
\textbf{Source problem code:} catalogue entry
\subsection{Definitions and mathematical statement}
Let $u:\R^n\to\R$ be continuous. The first property in Rubel's question is
\[
\exists\text{ a homeomorphism }\phi:\R^n\to\R^n\quad
\exists v\in C^2(\R^n;\R),\qquad
\Delta v=0,\quad v(x)=u(\phi(x))\ \text{for all }x.
\]
Here $\Delta v=\sum_{j=1}^n\partial_{x_j}^2v$. No differentiability of $\phi$ is imposed. The question seeks a necessary and sufficient topological condition for this harmonic reparametrization.

The proposed second property says that there exist coordinate functions $\mu_1,\ldots,\mu_n$ such that
\[
U=(u,\mu_1,\ldots,\mu_n)
\]
is a light open map ``of $\R^n$ into $\R^n$''. A continuous map is open when it takes open subsets to open subsets of the specified target, and light when every point inverse is totally disconnected. The source calls the $\mu_j$ mappings without explicitly stating their continuity; continuity of $U$ is the usual light-open convention.

As printed, $U$ has $n+1$ scalar coordinates and hence takes values in $\R^{n+1}$, not $\R^n$. Therefore the proposed property is not well typed. The usual coordinate-count repair would use $n-1$ auxiliary real functions, giving $(u,\mu_1,\ldots,\mu_{n-1}):\R^n\to\R^n$, but this is identified as a possible correction rather than silently substituted. Changing the target to $\R^{n+1}$ is a different correction and not equivalent by notation alone.

There are further missing qualifications even after a coordinate-count repair. Constant $u$ satisfies the harmonic-reparametrization property, whereas a nonempty map with a constant first coordinate has image in a hyperplane and cannot be open into $\R^n$. Thus a nonconstant restriction or separate treatment of constants is required for the intended equivalence. The word ``into'' does not require surjectivity; this matters for any global claim. Local light-open factorization and a global harmonic function on all of $\R^n$ must also be distinguished.

The source attributes the case $n=2$ to Stoilow, but its citation is missing. The exact planar theorem, global conformal-type assumptions, and treatment of the range must be checked rather than inferred from the incomplete attribution. The supplied triage labels the item open as of the 2018 edition and reports no recent resolution; no independent status assessment is supplied here.
\subsection{Short English statement}
When can a continuous function on Euclidean space become harmonic after a homeomorphism of the whole space? The proposed light-open completion criterion has a coordinate-count error and omits a necessary treatment of constant functions. Its claimed planar Stoilow case also needs the missing citation and global hypotheses before the equivalence can be stated precisely.
\subsection{Sources}
\begin{itemize}
\item Local source registry: \url{data/unsolvedmath/problems.json}, permanent UnsolvedMath ID $2303020$, source code catalogue entry, including the supplied OPEN-TRIAGE assessment. Dataset attribution: Ulam AI and UnsolvedMath Contributors; original-source attribution and rights are retained.
\item Hayman, \emph{Research Problems in Function Theory}, source-list date 2018, Problem 3.20, proposed by L. A. Rubel and communicated by D. A. Brannan: \url{https://arxiv.org/abs/1809.07200}. The metadata reports source-TeX extraction and July 29, 2026 access, with current-status and rights review flags.
\item The statement attributes its $n=2$ case to Stoilow, but the phrase ``which is in , for example'' has no reference in the extract. No exact planar theorem or separate authoritative citation is supplied. No linked source was inspected for this draft.
\end{itemize}
\end{document}
